%!TEX root = ../thesis.tex
\setlength{\parskip}{0pt}
\chapter{基于多表的增量式跨模态哈希检索算法}
\label{chap:mih}
\section{引言}
上一章围绕存在概念漂移的动态数据环境下的哈希检索问题展开研究，针对单模态图像数据，提出了一种基于概念保留的动态半监督哈希算法SCPH。该方法在无需对历史数据的哈希编码进行更新的前提下，有效增强了哈希模型对数据分布变化的适应能力，相关实验结果验证了所提出方法的有效性与可行性。然而，在实际应用场景中，检索任务往往涉及多种模态数据的联合建模与处理，不同模态在特征表示形式和语义表达层面存在显著差异，使得模型在动态数据环境下面临更加复杂的语义一致性保持与跨模态对齐问题。

在动态多模态数据环境中，不同模态的数据通常以数据流的形式持续到达，数据分布随时间不断演化，新语义类别的引入以及既有类别分布的变化使跨模态样本之间的相似性关系呈现出明显的不稳定性。当概念漂移现象发生时，基于历史数据学习得到的跨模态哈希模型往往难以有效适应当前的数据分布，从而导致检索性能显著下降。此外，现有多数动态跨模态哈希方法通常依赖单一的哈希表，并在新数据到来时对模型进行频繁更新。在概念漂移较为显著的情况下，新旧数据之间的语义差异不断累积。在增量学习过程中，这类方法易产生灾难性遗忘现象，使模型在对新数据进行适配的同时，对历史跨模态相似性结构的保持能力显著下降，进而影响检索性能的长期稳定性与可靠性。

针对上述问题，本章提出一种基于多表的增量跨模态哈希方法（Multi-table based Incremental Hashing，MIH），用于解决概念漂移环境下的跨模态哈希检索问题。MIH方法通过构建多哈希表系统，在不同时间步分别学习并保留跨模态相似性信息，以缓解概念漂移与灾难性遗忘对模型性能的影响；同时，根据各哈希表在最新到达的多模态数据上的表现自适应地分配权重，并在检索阶段引入查询自适应的加权策略，从而更充分地利用不同哈希函数对查询模态数据的判别能力，生成更加有效的汉明距离度量，提升跨模态检索性能。

\section{MIH算法介绍}
MIH主要聚焦于图像与文本之间的跨模态检索任务，该算法具有良好的通用性和可扩展性，可自然推广至更多模态的数据场景。因此，不失一般性的，我们假设多模态数据以图像-文本对的形式随时间的推移连续到达，在 $t$ 时刻，新到来的图像-文本数据块记为 $\mathbf{X}_t=\{\mathbf{X}_1^t, \mathbf{X}_2^t\}$，其中 $\mathbf{X}_1^t \in \mathbb{R}^{d_1 \times n}$ 和 $\mathbf{X}_2^t \in \mathbb{R}^{d_2 \times n}$ 分别表示图像模态和文本模态的特征表示，$n$ 表示当前时刻 $t$ 到达的数据块中所包含的样本数量，$d_1$ 和 $d_2$ 分别表示图像与文本模态的特征维度。此外，$\mathbf{L}_t \in \{0,1\}^{c_t \times n}$ 表示 $t$ 时刻新到达数据所对应的标签矩阵，其中，$c_t$ 表示截止至 $t$ 时刻所观测到的标签类别总数。当 $L_{ij}=1$ 时，表示第 $j$ 个样本属于第 $i$ 个类别，若为 $0$，则表示第 $j$ 个样本不属于第 $i$ 个类别。在 $t$ 时刻，MIH待学习的模态共享的哈希编码表示为 $\mathbf{B}_t \in \{-1,1\}^{r \times n}$，其中 $r$ 表示哈希编码的长度。$\mathbf{B} \in \{-1,1\}^{r \times n_{t-1}}$ 表示在 $t$ 时刻之前学习到的模态共享的哈希编码。$\mathbf{W}_1^t \in \mathbb{R}^{r \times d_1}$ 和 $\mathbf{W}_2^t \in \mathbb{R}^{r \times d_2}$ 分别表示图像模态的哈希映射函数和文本模态的哈希映射函数。

在 $t$ 时刻，MIH 构建一个包含 $K$ 个跨模态哈希表的多表跨模态哈希系统。每一个跨模态哈希表由一组图像模态与文本模态的哈希函数以及对应的表权重组成。该 $K$ 个哈希表分别基于不同时刻到达的数据进行训练，从而能够在时间维度上对不同阶段的数据语义结构进行建模，并持续保留各阶段所学习到的跨模态相似性信息。为适应当前数据分布的变化，MIH 根据最新到达的数据对各跨模态哈希表进行性能评估，并据此分配相应的表权重，使得系统在检索过程中能够更加有效地利用对当前数据更具判别能力的哈希表。在 $t$ 时刻，MIH 可表示为如下形式：
\begin{equation}
  \Psi_t = \{(v_1^{tk}, \mathbf{W}_1^{tk}), (v_2^{tk}, \mathbf{W}_2^{tk})\}, \quad k \in \{1,2,\dots,K\}
\end{equation}
其中，$v_1^t$ 和 $v_2^t$ 分别表示 $t$ 时刻下，图像模态和文本模态哈希函数的权重。当最新数据块 $\mathbf{D}_t$ 到达时，MIH 基于 $\mathbf{D}_t$ 采用 FCMH 算法学习当前时间步的数据哈希编码 $\mathbf{B}_t$，以及对应的图像模态与文本模态哈希函数 $\mathbf{W}_1^t$ 和 $\mathbf{W}_2^t$。新生成的哈希编码 $\mathbf{B}_t$ 与此前各时间步已学习得到的历史哈希编码共同构成更新后的检索库 $\mathbf{B}$，用于后续跨模态检索任务。

在 $t$ 时刻，MIH中的多表跨模态哈希系统暂时包含 $K+1$ 组哈希表结构，其中包括 $t-1$ 时刻中 $\Psi_{t-1}$ 保留的 $K$ 组跨模态哈希表，以及基于当前数据块训练得到的哈希表 $\Psi_0$。对于这 $K+1$ 组跨模态哈希表，MIH使用最新到达的数据对齐进行评估，并相应更新各自的权重。为了控制存储开销并保证检索效率，MIH在 $t$ 时刻仅保留 $K$ 组跨模态哈希函数以及对应的权重。因此，在完成权重评估之后，MIH将根据哈希表的权重，从当前的 $K+1$ 组哈希表中移除权重最低的一组哈希表，维持固定规模的多表结构。

在检索阶段，考虑到查询样本相对于哈希超平面的距离差异会影响其哈希编码的判别能力，MIH针对跨模态哈希函数引入了一种查询自适应加权机制，以增强查询样本哈希表示的区分性。具体而言，在计算相似度时，MIH综合多表跨模态哈希系统中各哈希表的权重以及与查询样本相关的自适应权重，对汉明距离进行加权计算，从而获得更加精确的相似度度量，提升整体跨模态检索性能。

\subsection{MIH中哈希表的训练}
在 MIH 算法框架中，多表跨模态哈希系统中的每一组哈希表均通过 FCMH\upcite{Wang2022FCMH} 算法进行训练。鉴于 FCMH 并非本文的创新点，本文仅对该算法进行简要介绍。FCMH 采用两阶段学习策略构建跨模态哈希模型：首先基于训练数据学习统一的跨模态哈希码；随后通过线性回归方式学习各模态对应的哈希函数。具体来说，FCMH利用标签全局相似性和局部相似性来指导哈希码的学习。首先，FCMH利用样本间的类别标签信息构建全局相似性关系，并将其嵌入到哈希码的学习过程中，可表示为如下：
\begin{equation}
  \begin{aligned}
    \min_{\mathbf{B}_t,\mathbf{P}} \quad
    & \alpha \|(\mathbf{B}_t)^{\top} \mathbf{B}_t - r \mathbf{S}_g\|_F^2 + \gamma \|\mathbf{P}\|_F^2 \\
    & + \beta \sum_{i,j=1}^{c_t} [\mathbf{C}_g]_{ij}
    \|[\mathbf{L}_t]_{i*} - [\mathbf{P}\mathbf{B}_t]_{j*}\|_2^2 \\
    \text{s.t.} \quad
    & \mathbf{B}_t \in \{-1,1\}^{r \times n}, \\
    & \mathbf{S}_g = 2\mathbf{L}_t^{\top}\mathbf{L}_t - \mathbf{1}\mathbf{1}^{\top}
  \end{aligned}
\end{equation}
其中，$\mathbf{S}_g \in [-1,1]$ 是全局成对相似性矩阵，$\mathbf{P}$ 是一个映射矩阵，$\alpha$、$\beta$、$\gamma$ 是平衡超参数，$\mathbf{L}_t^{\top}$ 是 $t$ 时刻 $\text{L}_2$ 范数列归一化的标签矩阵，它的第 $j$ 列定义为 $[\mathbf{L}_t]_{*j} = [\mathbf{L}_t]_{*j} / \|[\mathbf{L}_t]_{*j}\|$，$\mathbf{C}_g \in \mathbb{R}^{c \times c}$ 是全局类别相关性矩阵，其定义如下：
\begin{equation}
  \mathbf{C}_g = \mathbf{L}_t\mathbf{L}_t^{\top}
\end{equation}
其中，$\mathbf{L}_t$ 是 $\text{L}_2$ 范数行归一化的标签矩阵，它的第 $i$ 行定义为 $[\mathbf{L}_t]_{i*} = [\mathbf{L}_t]_{i*} / \|[\mathbf{L}_t]_{i*}\|$，$[\mathbf{C}_g]_{ij} \in [0,1]$ 的值越大，表明 $i$ 类别和 $j$ 类别同时出现的频率越高。也就是说，如果一个样本属于 $i$ 类，那么它属于 $j$ 类的可能性也很大。

尽管全局相似性矩阵 $\mathbf{S}_g$ 和 $\mathbf{C}_g$ 能够保留在粗粒度层面保留样本间的语义关系，但依赖类别统计信息，会忽略局部领域结构中的细粒度信息。尤其在多标签场景下，具有相同或相关标签的样本在特征空间中仍可能存在差异。因此，FCMH在全局语义建模基础上引入局部相似性约束，以保留样本的局部结构信息并提升细粒度语义表达能力。

具体而言，首先，假设将第 $m$ 模态的训练集划分为 $p$ 组，即 $\mathbf{X}_t^{(m)} = [\mathbf{X}_1^{(m)}, \mathbf{X}_2^{(m)}, \dots, \mathbf{X}_p^{(m)}]$，其中 $\mathbf{X}_q^{(m)} \in \mathbb{R}^{d_m \times n_q} (q \in \{1,2,\dots,p\})$。这些分组可通过先验知识或聚类生成。$\mathbf{L}_q$、$\mathbf{S}_q^l$ 和 $\mathbf{B}_q$ 分别是 $\mathbf{L}$、$\mathbf{S}$ 和 $\mathbf{B}$ 的第 $q$ 子集。与全局相似性矩阵 $\mathbf{S}_g$ 和全局类相关矩阵 $\mathbf{C}_g$ 的定义类似，第 $q$ 组的局部成对相似性矩阵 $\mathbf{S}_q^l$ 与局部类相关性矩阵$\mathbf{C}_q^l$可定义为如下：
\begin{align}
  \mathbf{S}_q^l &= 2\mathbf{L}_q^{\top}\mathbf{L}_q - \mathbf{1}\mathbf{1}^{\top}
  \label{eq:mih-local-sim} \\
  \mathbf{C}_q^l &= \mathbf{L}_q\mathbf{L}_q^{\top}
  \label{eq:mih-local-corr}
\end{align}
因此，类似于公式4-2，将局部相似性嵌入到哈希码中的目标函数可定义为如下：
\begin{equation}
  \begin{aligned}
    \min_{\mathbf{B}_t,\mathbf{P}} \quad
    & \sum_{q=1}^{m} \alpha_q \|\mathbf{B}_q^{\top} \mathbf{B}_q - r \mathbf{S}_q^l\|_F^2 + \gamma \|\mathbf{P}\|_F^2 \\
    & + \sum_{q=1}^{m} \beta_q \sum_{i,j=1}^{c_t} [\mathbf{C}_q^l]_{ij}
    \|[\mathbf{L}_q]_{i*} - [\mathbf{P}\mathbf{B}_q]_{j*}\|_2^2 \\
    \text{s.t.} \quad
    & \mathbf{B}_t \in \{-1,1\}^{r \times n}
  \end{aligned}
\end{equation}
结合公式4-2和4-6，FCMH的总体目标函数如下：
\begin{equation}
  \begin{aligned}
    \min_{\mathbf{B}_t,\mathbf{P}} \quad
    & \alpha_1 \|(\mathbf{B}_t)^{\top} \mathbf{B}_t - r \mathbf{S}_g\|_F^2
    + \alpha_2 \sum_{q=1}^{m} \|\mathbf{B}_q^{\top} \mathbf{B}_q - r \mathbf{S}_q^l\|_F^2 \\
    & + \gamma \|\mathbf{P}\|_F^2
    + \beta_1 \sum_{i,j=1}^{c_t} [\mathbf{C}_g]_{ij}
    \|[\mathbf{L}_t]_{i*} - [\mathbf{P}\mathbf{B}_t]_{j*}\|_2^2 \\
    & + \beta_2 \sum_{q=1}^{m} \sum_{i,j=1}^{c_t} [\mathbf{C}_q^l]_{ij}
    \|[\mathbf{L}_q]_{i*} - [\mathbf{P}\mathbf{B}_q]_{j*}\|_2^2 \\
    \text{s.t.} \quad
    & \mathbf{B}_t \in \{-1,1\}^{r \times n}
  \end{aligned}
  \label{eq:mih-fcmh}
\end{equation}
通过离散迭代优化公式4-7，可以得到最优的模态共享哈希码 $\mathbf{B}_t$。随后，通过线性回归方法分别为不同模态学习对应的模态哈希映射函数 $\mathbf{W}$，从而实现从原始特征空间到汉明空间的映射。

\subsection{哈希表权重的计算方法}
在 MIH 框架中，哈希表中各模态对应的哈希函数权重用于评估新到达数据的跨模态相似性保持能力。具体而言，哈希表的整体权重用于度量各哈希表的性能表现，并据此权重最低的哈希表进行淘汰；各模态哈希函数的权重则用于在跨模态检索过程中对相似性结果加权汉明距离计算。此外，每个模态哈希函数权重的均值被定义为该哈希表的最终权重。鉴于本节主要讨论同一时间步 $t = T$ 下的情形，为简化表示，以下公式中省略上标 $t$。

当最新的数据块 $\mathbf{X}=\{\mathbf{X}_1, \mathbf{X}_2\}$ 到达时，MIH使用公式4-7学习到模态共享的哈希码 $\mathbf{B} \in \{-1,1\}^{r \times n}$ 以及对应的图像和文本模态的哈希映射函数 $\mathbf{W}_1$ 和 $\mathbf{W}_2$。基于模态映射函数，我们可以为不同模态的数据生成其对应的模态共享的哈希编码，公式如下：
\begin{equation}
  \mathbf{B}_i' = \operatorname{sign}(\mathbf{W}_i \mathbf{X}_i)
\end{equation}
其中，$i \in \{1,2\}$ 分别代表图像模态和文本模态。

在跨模态检索任务中，通常首先利用查询样本所属模态对应的哈希映射函数，将其投影至共享的汉明空间中，并通过计算其与检索库样本之间的汉明距离来度量相似性。一般而言，样本间语义相似度越高，其对应的汉明距离越小；反之，则汉明距离越大。因此，MIH通过度量不同模态生成的哈希码与直接学习得到的哈希码之间的一致性程度，来衡量哈希函数对于跨模态语义关系的保留能力。具体而言，对于新到达的数据，首先利用公式4-8生成各个模态对应的哈希码 $\mathbf{B}_i'$，随后与已学习得到的模态共享哈希码 $\mathbf{B}$ 进行汉明距离计算，以反映不同模态哈希表示在语义空间中的一致性，具体公式如下：
\begin{equation}
  \mathbf{Ham}_i = \frac{1}{2} (r \mathbf{1}_{n \times n} - (\mathbf{B}_i')^{\top} \mathbf{B})
\label{eq:mih-ham}
\end{equation}
其中 $\mathbf{1}_{n \times n}$ 表示全 1 矩阵。在MIH中，通过公式4-7学习得到的哈希码是作为新数据在汉明空间中的语义表示基准，因此，通过模态隐射函数 $\mathbf{W}$ 得到的哈希码应尽可能的逼近 $\mathbf{B}$，从而实现跨模态语义对齐。基于这种跨模态一致性，可以衡量各模态哈希函数对于语义信息的保留能力，并据此计算各模态哈希函数的权重，公式如下：
\begin{equation}
  v_i = \sum_{\mathbf{S}_{jk}=1} [\mathbf{Ham}_i]_{jk} - \sum_{\mathbf{S}_{jk}=-1} [\mathbf{Ham}_i]_{jk} = \operatorname{sum}(\mathbf{S} \circ \mathbf{Ham}_i)
\label{eq:mih-modal-weight}
\end{equation}
其中，$\mathbf{S} \circ \mathbf{Ham}_i$ 表示矩阵 $\mathbf{S}$ 和矩阵 $\mathbf{Ham}_i$ 之间的逐元素积，$\operatorname{sum}(\cdot)$ 表示对矩阵内所有元素进行求和，$v_i$ 表示第 $i$ 个模态的哈希函数的权重。$\mathbf{S}$ 是成对相似性矩阵，其中 $\mathbf{S}_{jk} \in \{-1,1\}$，当 $\mathbf{S}_{jk}=1$ 时，表示第 $j$ 个样本与第 $k$ 个样本至少共享一个语义标签，反之亦然。

最终，将模态间一致性权重（即各模态哈希函数权重）归一化至 $[0, 1]$ 范围，具体公式如下：
\begin{equation}
  v_i = \frac{v_i - \min(v_i)}{\max(v_i) - \min(v_i)}
\label{eq:mih-modal-weight-norm}
\end{equation}
其中，$\max(v_i)$ 和 $\min(v_i)$ 分别为 $v_i$ 的可取的最大值和最小值。通过理论分析可知，$v_i$ 的理论最大值和最小值分别为 $r \mathbf{S}_1$ 和 $-r \mathbf{S}_{-1}$，其中 $\mathbf{S}_1$ 为相似样本的数量，$\mathbf{S}_{-1}$ 为不相似样本的数量。将各模态对应的哈希函数权重进行求和并取均值，作为该哈希表的整体权重指标，用于后续淘汰表过程中淘汰权值最小的哈希表。

\subsection{基于查询自适应权重与哈希表权重的加权检索方法}
在动态数据环境中，由于数据的分布可能会发生变化，不同哈希表在当前阶段的检索性能可能存在差异。因此，MIH基于最新时刻到达的数据为哈希表赋予相应的权重。此外，在多表跨模态哈希系统中，不同的哈希函数对应的超平面对查询样本的区分程度不同，从而导致其对检索结果的贡献存在差异。

因此，本文进一步计算查询样本与对应模态的哈希函数超平面之间的距离，并将该距离引入作为模态相关的查询自适应权重，以更细粒度的表征各哈希比特位在当前查询样本下的贡献差异，进一步提升哈希码各比特位的判别能力。具体而言，设第 $i$ 个模态的查询样本为 $\mathbf{x}_i^q$，第 $i$ 模态的哈希函数中的一个哈希函数为 $\mathbf{w}_i$，则查询样本与该哈希超平面的欧式距离可表示为如下：
\begin{equation}
  D(\mathbf{x}_i^q, \mathbf{w}_i) = \frac{\mathbf{x}_i^q \cdot \mathbf{w}_i}{\|\mathbf{w}_i\|}
\label{eq:mih-query-distance}
\end{equation}
其中，$D(\cdot,\cdot)$ 表示查询样本与对应模态的哈希超平面之间的欧式距离，$\|\cdot\|$ 表示向量的范数，$\mathbf{x}_i^q \cdot \mathbf{w}_i$ 表示查询样本与哈希超平面向量的点积。将公式4-12归一化至 $[0, 1]$ 区间内，作为最终的模态特定的查询自适应权重，归一化公式如下：
\begin{equation}
  \alpha_i(\mathbf{x}_i^q, \mathbf{w}_i) = \frac{D(\mathbf{x}_i^q, \mathbf{w}_i) - D_{\min}}{D_{\max} - D_{\min}}
\label{eq:mih-query-weight}
\end{equation}
其中，$\alpha_i(\mathbf{x}_i^q, \cdot)$ 表示第 $i$ 个模态的查询样本的查询自适应权重，$D_{\max}$ 和 $D_{\min}$ 分别表示 $D(\mathbf{x}_i^q, \mathbf{w}_i)$ 的最大值和最小值。基于模态间一致性权重 $v_i$ 以及查询自适应权重 $\alpha_i$，对于第 $i$ 模态的查询样本 $\mathbf{x}_i^q$ 与待检索集 $\mathbf{B}$ 之间的最终加权汉明距离可通过如下式子计算：
\begin{equation}
  H_i(\mathbf{x}_i^q, \mathbf{B}) = \sum_{k=1}^{K} v_{ik} \sum_{j=1}^{r} \alpha_i(\mathbf{x}_i^q, \mathbf{w}_{ij}) \left( h_k^j(\mathbf{x}_i^q) \oplus h_j(\mathbf{x}) \right)
\label{eq:mih-weighted-hamming}
\end{equation}
其中，$h_k^j(\mathbf{x}_i^q)$ 表示由第 $k$ 个哈希表为查询样本的生成的第 $j$ 位哈希码，$h_j(\mathbf{x})$ 表示所有待检索集的第 $j$ 位模态共享哈希码，$\oplus$ 表示用于计算汉明距离的异或操作。

综合上述过程，MIH 在每个时间步的核心流程可概括为：利用当前数据块增量训练新哈希表，基于跨模态一致性计算并更新哈希表权重；当哈希表数量超过$K$时，淘汰权重最低的哈希表；在检索阶段，联合哈希表权重与查询自适应权重计算最终加权汉明距离并完成排序。其伪代码流程如算法\ref{alg:mih-update-retrieval} 所示。

\vspace{0.3\baselineskip}
\begingroup
\setlength{\intextsep}{4pt}
\begin{algorithm}[H]
  \caption{时间步 $t$ 的 MIH 训练与检索流程}
  \label{alg:mih-update-retrieval}
  \begin{algorithmic}[1]
    \REQUIRE 当前数据块 $\mathcal{D}_t$，最大哈希表数量 $K$，当前哈希表集合 $\mathcal{T}_{t-1}$
    \ENSURE 更新后的哈希表集合 $\mathcal{T}_{t}$，查询样本的检索排序结果
    \STATE 基于 $\mathcal{D}_t$ 和 式\eqref{eq:mih-fcmh} 训练新哈希表 $T_t$，得到模态共享哈希码 $\mathbf{B}$ 及各模态哈希函数
    \STATE 根据式\eqref{eq:mih-ham} 计算各模态汉明距离矩阵 $\mathbf{Ham}_i$
    \STATE 根据式\eqref{eq:mih-modal-weight} 和式\eqref{eq:mih-modal-weight-norm} 计算新表权重
    \STATE 将 $T_t$ 加入哈希表集合：$\mathcal{T}_{t} \leftarrow \mathcal{T}_{t-1} \cup \{T_t\}$
    \IF{$|\mathcal{T}_{t}| > K$}
      \STATE 淘汰权重最小的哈希表
    \ENDIF
    \FOR{每个查询样本 $\mathbf{x}_i^q$}
      \STATE 根据式\eqref{eq:mih-query-distance} 和式\eqref{eq:mih-query-weight} 计算查询自适应权重
      \STATE 根据式\eqref{eq:mih-weighted-hamming} 计算加权汉明距离
      \STATE 按距离升序排序并返回检索结果
    \ENDFOR
  \end{algorithmic}
\end{algorithm}
\endgroup

\section{实验与分析}
\subsection{实验基本设置}
在实验部分，为了系统评估MIH方法的检索性能，本文选取两个常用的跨模态检索基准数据集——MIRFlickr\upcite{Huiskes2008MIRFlickr} 数据集和 MSCOCO 数据集进行实验验证\upcite{Lin2014MSCOCO}。

MIRFlickr 数据集包含 25,000 对从 Flickr 网站收集的图文样本，每一对图文数据至少关联 24 个预定义语义类别中的一个标签。为保证实验数据的有效性与特征表达质量，在数据预处理阶段，首先删除出现频次少于 20 次的文本标签，以减少稀疏标签对模型训练带来的干扰；随后，剔除缺少文本标签或人工标注信息的样本。经过上述处理后，最终保留 20,015 对有效图文样本用于后续实验。在特征表示方面，每幅图像采用 512 维 GIST 特征进行表示，文本数据则使用 1,386 维词袋模型（Bag-of-Words，BoW）特征进行编码，从而构建对应的跨模态特征空间。

MSCOCO 数据集包含 122,218 对带有语义标注的图文样本，这些样本均来源于 Flickr 网站。该数据集共划分为 80 个子类别，并进一步归属于 12 个超类别。每一对图文样本至少对应 80 个子类别中的一个语义标签，从而为跨模态语义匹配任务提供明确的监督信息。在特征表示方面，图像数据采用 512 维 GIST 特征进行表征；文本数据首先基于原始 1,000 维词袋模型（Bag-of-Words，BoW）特征进行主成分分析（Principal Component Analysis，PCA）降维处理，最终得到 150 维特征向量用于表示文本信息。

为了有效的评估MIH方法在概念漂移场景下的检索性能，本文模拟了两种存在概念漂移的动态数据环境，分别是新类出现和数据分布随时间变化两种典型的概念漂移场景。每个场景均设置了15个时间批次。首先，使用MIRFlickr 数据集模拟新类出现的动态场景。在该设置中，每个时间步均引入一个新的数据块，其中训练集包含 1,000 个样本，对应测试集包含 200 个样本。在 $t = 1$至 $t = 5$ 阶段，数据块中的样本均从随机选取的 15 个类别中抽取。自 $t = 6$ 起，在原有类别基础上逐步引入剩余的 9 个类别，并从原始类别与新引入类别中随机采样生成新的数据块，从而模拟类别空间逐步扩展的过程，体现“新类出现”这一概念漂移场景。其次，使用MSCOCO 数据集构建数据分布改变的漂移场景。在该场景中，每个时间步的数据块包含 2,000 个训练样本及 500 个测试样本。为模拟分布变化，在每个超类别内部随机划分子类别，其中部分子类别被设定为新出现子类别，其余视为旧子类别。具体而言，从各超类别的全部子类别中选取 70\% 作为新出现子类别。在 $t = 1$ 至 $t = 5$ 阶段，各时间步的数据块及对应测试集均仅从旧子类别样本中随机抽取；自 $t = 6$ 起，各超类别内部逐步引入新子类别样本，导致整体数据分布结构发生变化，从而模拟数据分布逐渐发生改变的过程。

为全面验证所提出 MIH 方法在动态环境下的有效性与泛化能力，本文选取了多种具有代表性的跨模态哈希方法作为对比，包括 FDDH-online\upcite{liu2021fddh}、OSCMFH\upcite{shu2023online}、ROH\upcite{jiang2023random}、ROHLSE\upcite{li2024robust}、OHPKL\upcite{shu2024online}、BATCH\upcite{Wang2021BATCH}、ALECH\upcite{9507286}、AMSH\upcite{10044981} 以及 FCMH\upcite{Wang2022FCMH}。这些方法涵盖离线方法、在线方法以及半监督在线方法，能够从不同学习范式角度对所提方法进行系统比较。其中，BATCH、ALECH、AMSH 和 FCMH 为离线跨模态哈希方法；OHPKL 为在线半监督方法；其余方法为监督式在线跨模态哈希方法。
在训练机制设计上，为保证实验公平性并突出在线方法在动态环境中的优势，离线跨模态哈希方法仅在第一个时间步进行训练，随后模型参数保持固定，不再进行参数更新，也不利用新到达的数据进行增量学习；新到达的数据由已训练完成的哈希模型进行编码。该设置能够检验离线模型在数据分布变化或类别扩展情况下的泛化能力。对于部分仅支持固定类别数量的在线方法，由于其理论上无法处理新类别出现的情形，在类别扩展场景中，每个时间步仅使用原始类别样本进行训练，从而避免引入额外的结构性偏差。在参数设置方面，最大哈希表数量 $K$设为 5，以平衡检索效率与表达能力。哈希编码长度设置为64。在每个时间步中，采用 FCMH 训练哈希表时的超参数严格遵循原论文设定，以确保对比结果的可复现性与公正性。

在任务层面，本文从双向跨模态检索角度对模型性能进行评估，包括图像到文本检索（I2T）与文本到图像检索（T2I）。这种双向评估能够全面反映模型在不同模态查询条件下的语义对齐能力。评价指标采用基于语义真实标签的 Top-100 mAP，以衡量模型在前 100 个检索结果中的整体排序质量。为降低随机性对结果的影响，所有实验均独立运行 10 次，并使用平均性能作为最终结果。
\subsection{实验结果分析}
在存在概念漂移现象的动态数据环境中，MIH与多种代表性对比方法的检索性能的对比结果如下。

\begin{figure}[htbp]
  \centering
  \includegraphics[width=0.8\textwidth]{MIH/MIRflickr/4-1}
  \caption{MIH和对比算法在MIRFlickr新类出现上的Top-100 mAP (I2T)}
  \label{fig:mirflickr-1}
\end{figure}

\begin{figure}[htbp]
  \centering
  \includegraphics[width=0.8\textwidth]{MIH/MIRflickr/4-2}
  \caption{MIH和对比算法在MIRFlickr新类出现上的Top-100 mAP (T2I)}
  \label{fig:mirflickr-2}
\end{figure}

图4-1和图4-2展示了在新类别逐步出现的动态数据场景下，各方法在 Top-100 mAP 指标上的检索性能变化情况。在I2T任务中，MIH在初始阶段的性能并非最优，但随着时间推进以及新类别的持续引入，其性能逐步提升，并最终达到与其他方法相当甚至更优的水平。这表明，在类别分布相对稳定的早期阶段，各方法均能够较好地拟合当前数据结构，因此性能差异有限；而当类别空间不断扩展、语义结构发生变化时，MIH 所采用的多哈希表机制能够并行维护不同时间阶段的语义信息，有效缓解新旧知识之间的冲突，从而逐步体现出性能优势。尤其在 T2I 任务中，MIH 在所有时间步均取得最优结果，并在新类别出现后优势更加明显，说明该方法在类别扩展情形下具有较强的适应能力。

\begin{figure}[htbp]
  \centering
  \includegraphics[width=0.8\textwidth]{MIH/MSCOCO/4-3}
  \caption{MIH和对比算法在MSCOCO数据分布改变上的Top-100 mAP (I2T)}
  \label{fig:mscoco-1}
\end{figure}

\begin{figure}[htbp]
  \centering
  \includegraphics[width=0.8\textwidth]{MIH/MSCOCO/4-4}
  \caption{MIH和对比算法在MSCOCO数据分布改变上的Top-100 mAP (I2T)}
  \label{fig:mscoco-2}
\end{figure}

图4-3和图4-4给出了数据分布发生改变场景下的 Top-100 mAP 检索结果。可以观察到，在I2T与T2I 两种检索任务中，MIH 在所有时间步均持续优于其他对比方法，且性能波动较小。与新类别出现情形不同，数据分布改变更强调语义结构的渐进性变化。在此情况下，MIH 通过语义一致性权重排序机制，根据各哈希函数在最新数据上的表现动态分配权重，从而提高当前更具判别能力的哈希函数在检索过程中的贡献度；同时，查询自适应加权策略增强了查询样本在哈希空间中的区分能力，从而提升汉明距离对语义相似样本的区分能力。因此，在数据分布持续变化的动态环境中，MIH 仍能够保持稳定而优越的检索性能。
对比方法中，BATCH、ALECH、AMSH 和 FCMH 属于静态跨模态哈希方法，其模型基于静态数据假设进行一次性训练，参数在后续阶段不再更新，因此在动态数据环境下难以适应数据分布变化，导致检索性能下降。其他在线跨模态哈希方法虽具备增量学习能力，但通常依赖单一模型结构进行持续更新，缺少用于保留历史语义信息或协调新旧知识关系的专门机制，并且忽略了概念漂移问题。当类别空间扩展或语义结构发生变化时，模型难以有效协调新旧知识之间的关系，因而出现性能退化。
相比之下，MIH 通过多哈希表并行维护不同时间阶段的语义表示，并结合基于最新数据表现的动态加权策略，实现对历史信息的保留与当前分布的自适应平衡。同时，查询自适应加权增强了哈希表示的判别能力，使得基于汉明距离的排序结果在分布变化条件下仍能保持较高的语义一致性。因此，在存在概念漂移的动态环境中，MIH表现出更好的稳定性与检索精度。

\subsection{消融实验}
为系统评估 MIH 中模态间一致性加权机制与查询自适应加权机制在检索阶段的作用，本文构建了三种变体模型，并与完整的 MIH 进行对比分析，分别记为 MIH1、MIH2 和 MIH3。具体而言，MIH1 在加权汉明距离计算过程中移除模态间一致性权重；MIH2 移除查询自适应权重；MIH3 同时移除上述两类加权机制。通过构造不同程度的简化模型，可以分别分析各权重机制对检索性能的影响。上述模型均在 MSCOCO 数据集上进行验证，并在文本到图像检索（T2I）任务下进行性能比较。各方法在 Top-100 mAP 指标上的检索结果如图4-5所示。

\begin{figure}[htbp]
  \centering
  \includegraphics[width=0.8\textwidth]{MIH/ablation/4-5}
  \caption{MIH及其变体在MSCOCO上的Top-100 mAP (T2I)}
  \label{fig:ablation}
\end{figure}

由图4-5可知，在初始阶段（$t = 1$ 至 $t = 5$），各方法均表现出相近且较高的检索精度，此时数据分布保持稳定，尚未出现概念漂移。自$t = 6$ 起，随着数据分布发生变化，各方法性能均出现不同程度下降，表明概念漂移对模型稳定性产生影响。在此过程中，完整的 MIH 始终优于其他三个消融变体。移除模态间一致性加权的 MIH1与移除查询自适应加权的MIH2 均表现出明显性能退化，说明两类加权机制分别在语义融合与查询判别能力增强方面发挥重要作用。进一步地，同时移除两类加权机制的 MIH3 性能最低，表明两种策略在应对分布变化时具有协同效应。因此，消融结果表明，模态间一致性权重与查询自适应权重均为 MIH 的关键组成部分，移除任一权重都会削弱模型对概念漂移的适应能力，进而导致检索性能的显著退化。

\section{本章小结}
本章针对存在概念漂移的动态数据环境下跨模态检索性能退化的问题，提出了一种基于多哈希表结构的增量跨模态哈希方法（Multi-table based Incremental Hashing, MIH）。MIH通过构建多哈希表系统，在时间维度上分别学习并保留不同阶段数据的语义结构，从而避免单一模型在持续更新过程中对历史语义信息的覆盖与遗忘，有效缓解了灾难性遗忘问题。

具体而言，MIH为每个时间步训练得到的跨模态哈希函数构建独立的哈希表，并根据其在最新到达数据上的跨模态语义保留能力计算相应权重，使在当前数据上表现更好的哈希表在检索阶段占据更高比重。同时，在检索阶段引入查询自适应权重，根据查询样本与各哈希函数超平面的距离，对不同哈希比特位进行加权，从而提升汉明距离在排序过程中的区分能力。通过多哈希表结构、表权重更新机制以及查询自适应加权策略的结合，MIH在保留历史语义信息的同时，提高了对当前数据分布的适应能力。在两类动态数据场景下的实验结果表明，MIH在各时间步均取得优于对比方法的检索性能。在新类别逐步引入和数据分布持续变化两种概念漂移情形下，MIH均能够保持较为稳定的性能表现。消融实验进一步说明，模态间一致性权重与查询自适应权重均对提升检索效果具有重要作用。

尽管MIH在当前动态数据设定下取得了稳定性能，但仍存在进一步拓展空间。未来工作将进一步扩展MIH框架，使其能够处理模态异步到达的流式多模态数据场景，即不同模态数据在不同时间步到达的情形。该设定更加贴近真实应用环境，有助于提升方法在复杂动态场景下的适用性。
